\subsection{Retained column isometries and positive traces at a fixed scale}
\label{sec:retained-column-positive-traces}

All inner products in this subsection are conjugate-linear in the
first variable. The comparison parameter is always the originally
prescribed positive number; no observed norm is substituted for it.

\begin{lemma}[The scalar-bearing retained column dictionary]
\label{lem:retained-column-isometry}
Let \(E,H,C_i,V_i\) be finite-dimensional complex Hermitian spaces.
Let \(J:E\to H\) be an antiunitary isometric embedding, \(V:E\to E\)
and \(U_i:V_i\to V_i\) unitary, and \(I_i:V_i\to C_i\) isometric.
Suppose the actual bilinear forms and period maps \(P_i:C_i\to H\)
satisfy
\[
 B_i(x,y)=a_i\langle JVy,P_iI_iU_ix\rangle,\qquad a_i\ne0.
 \tag{RCT1}
\]
There is a linear coisometry \(C:H\to E^\vee\) with
\[
 C(h)(y)=\langle JVy,h\rangle,\qquad
 CC^*=1,\qquad C^*C=\Pi_{J(E)},\qquad
 B_i=a_iCP_iI_iU_i.
\]
Here \(E^\vee\) is the complex bilinear dual with its dual norm.
Let \(j_i:V_i^0\to V_i\) be isometries and \(\Pi_0\) an
orthogonal projection on \(E^\vee\). Put
\(\Pi=C^*\Pi_0C\), \(J_i=I_iU_ij_i\), and fix the actual
normalizers \(c_i>0\). The retained maps
\(T_i=c_i^{-1/2}\Pi_0B_ij_i\) then have exact Gram blocks
\[
 T_i^*T_j=
 \frac{\bar a_i a_j}{\sqrt{c_ic_j}}\,
                 J_i^*P_i^*\Pi P_jJ_j
 \tag{RCT2}
\]
and target Gram
\[
 \sum_iT_iT_i^*
 =C\left(\sum_i\frac{|a_i|^2}{c_i}
                    \Pi P_iJ_iJ_i^*P_i^*\Pi\right)C^*.
 \tag{RCT3}
\]
\end{lemma}
\begin{proof}
The two conjugations in \(\langle JVy,h\rangle\) make it
linear in \(y\), and it is linear in \(h\) by convention.
Its dual norm is the norm of the orthogonal projection of
\(h\) onto \(J(E)\). Thus \(C\) restricts to a linear
isometry from \(J(E)\) onto \(E^\vee\), proving the two
coisometry identities. In particular \(\Pi\) is an orthogonal
projection with range in \(J(E)\), and
\(\Pi_0C=C\Pi\). Now (RCT1) gives
\[
 T_i=(a_i/\sqrt{c_i})\,C\Pi P_iJ_i.
\]
Taking adjoints and using \(C^*C\Pi=\Pi\) proves
(RCT2)--(RCT3).
\end{proof}

The phases of \(a_i\) can be absorbed in unitary source
coordinates. Their moduli cannot: the weights \(|a_i|^2/c_i\)
remain in the target Gram and in the diagonal source blocks.
A nonzero comparison scalar alone is not an isometry for a
preassigned arithmetic metric. Thus (RCT1), for the actual
theta and period maps, is an explicit application hypothesis.

\paragraph{The finite projector at the actual row evaluation.}
The finite part of the retained-row comparison has an exact
normalization. We spell it out to separate it from the remaining
onto-row and theta-period identities. Let \(K_f\) be the original
finite compact group, let \(K(g)\) be its normal congruence subgroup,
and put
\[
 \Gamma=K_f/K(g),\qquad v_g=\mu_f(K(g))\in(0,\infty).
\]
Write \(e_{\rm an}=\sum_{\gamma\in\Gamma}e_\gamma\gamma\)
for the original finite idempotent. The analytic operator is
\[
 G_{\rm an}f=\int_{K_f}
       \frac{e_{\bar k}}{v_g}R(k)f\,d\mu_f(k),
 \qquad R(k)f(x)=f(xk).
                                                        \tag{RP1}
\]
Thus its kernel is divided by the volume of the congruence
subgroup, not by \(|\Gamma|\) or by a compact-type dimension.

Let \(M\) be the prescribed automorphic multiplicity space,
with its finite action \(\rho\), and let
\(\sigma:\Gamma\xrightarrow{\sim}\Gamma_{\rm TC}\)
be the original identification of the two finite quotient groups.
For the prescribed common lowest vector \(v_0\), group
identification \(\iota\), and finite placement \(t\), the actual
translated evaluation is
\[
 E_t(T)(x)=T(v_0)(\iota^{-1}(x)t).
                                                        \tag{RP2}
\]
In the following lemma its values are assumed to belong to the
original \(L^2\) space. The two finite pins mean explicitly that
\[
 R(k)E_t(T)=E_t\bigl(\rho(\sigma(\bar k))T\bigr)
 \quad(k\in K_f),\qquad
 e_{\rm TC}(\sigma\gamma)=\overline{e_\gamma}.
                                                        \tag{RP3}
\]
The first equality is in \(L^2\), and the second uses the original
coefficient embedding. They are the source's level-action and
coefficient-matching clauses. Let \(\mathfrak c\) denote complex conjugation
of functions, an antiunitary map on this \(L^2\) space.

\begin{lemma}[Exact finite projector transport]
\label{lem:retained-finite-projector-transport}
Under (RP1)--(RP3), one has
\[
       \mathfrak c G_{\rm an} \mathfrak c\,E_t(T)=E_t(e_{\rm TC}T)
                                                        \tag{RP4}
\]
for every \(T\in M\), with no additional scalar. Consequently
\[
 E_t(e_{\rm TC}M)=\mathfrak c G_{\rm an}\mathfrak c\bigl(E_t(M)\bigr).
                                                        \tag{RP5}
\]
If the actual evaluation is additionally onto the prescribed
automorphic row space and satisfies its original norm identity
\(\|E_tT\|^2=c_{\rm ev}\|T\|_M^2\), with
\(c_{\rm ev}>0\), its normalization by
\(c_{\rm ev}^{-1/2}\) identifies these two retained spaces
isometrically.
\end{lemma}

\begin{proof}
Choose a representative \(k_\gamma\) of each reduction class.
The first pin implies that \(E_t(T)\) is right \(K(g)\)-invariant
and that its translate is constant, as an \(L^2\) vector, on
each fibre \(k_\gamma K(g)\). Every fibre has Haar measure
\(v_g\). Integrating (RP1) fibre by fibre therefore gives
\[
 G_{\rm an}f=\sum_{\gamma\in\Gamma}
                      e_\gamma R(k_\gamma)f
\]
for \(f\) in the evaluation image, and also for its complex
conjugate. The denominator \(v_g\) cancels the actual fibre
volume. Literal right translation commutes with complex
conjugation. Hence (RP3) gives
\[
 \begin{split}
 \mathfrak c G_{\rm an} \mathfrak c E_t(T)
   &=\sum_\gamma\overline{e_\gamma}R(k_\gamma)E_t(T)\\
   &=E_t\left(\sum_\gamma
           e_{\rm TC}(\sigma\gamma)\rho(\sigma\gamma)T\right)
     =E_t(e_{\rm TC}T).
 \end{split}
\]
The expression is independent of the chosen representatives.
Taking images proves (RP5). The additional positive norm identity
makes the normalized evaluation an isometry, and the additionally
assumed onto statement makes its restriction onto the retained
target. These last two properties have not been inferred from the
finite projector computation.
\end{proof}

\paragraph{The remaining common row map and its scalar.}
The normalized finite theta dictionary identifies the algebraic
terminal class pairing with its complex evaluation. To identify
that pairing with the literal Poincare columns, the missing
comparison is the same one in the source's \(\mathrm{RowTieT}\)
and \(\mathrm{MIid}\): a single semilinear row map \(e_R\)
must satisfy
\[
 \Theta_{\mu_0}(j_R e_Vy)=\lambda\,
              e_R(\iota_{\mu_0}c_Fy),\qquad \lambda\ne0,
                                                        \tag{RP6}
\]
on the original row quotient, be onto the specified good row
space, and have its original common norm ratio. Its translated
evaluation must be the same row realization used in (RP4).
The anchored see-saw identity must then hold with that map,
for every word and the actual columns.

For precision, write \(\epsilon\) for the fixed identity or
complex-conjugation coefficient transport. In the source's joined
identity let \(d.a_u\ne0\) be the anchored analytic scalar,
\(\lambda\) the scalar in (RP6), and \(\kappa_u\ne0\) the
pairing conversion between the same row evaluation and the
literal Poincare realization. Conjugating the row and column
norm-preserving maps through their prescribed similarities gives
unitaries \(V,U_u\); their norm ratios cancel in this conjugation.
The resulting scalar in the bilinear dictionary is exactly
\[
 a_u=\epsilon(d.a_u\lambda)\kappa_u,\qquad
 B_u(x,y)=a_u\langle JV y,\,P_u I_u U_u x\rangle.
                                                        \tag{RP7}
\]
Here the Hilbert product is conjugate-linear in its first argument,
\(J\) is the original antiunitary row inclusion, and \(I_u\)
is the original source inclusion. The formula follows first on
the algebraic coordinate bases from the base-tuple see-saw
identity, (RP6), and the pairing conversion; bilinearity extends
it to all complex arguments. With
\(\widehat P_u=c_{0X,u}^{-1/2}P_u\), the scalar in (RP7)
becomes \(a_u\sqrt{c_{0X,u}}\). Neither (RP4) nor the finite
theta dictionary proves that its modulus is one. The actual
onto-row and anchored identities in (RP6)--(RP7) remain to be
supplied; the finite-projector lemma does not assume their
conclusion under a change of notation.

\paragraph{The reference row identity and its exact evaluation factor.}
There are two further exact consequences of the original row pins.
They do not require choosing a new row space or a new normalizer.
The hypotheses in the following statement refer to the prescribed
complex theta model and Petersson realization; their construction
for the same moving family remains an application obligation.

\begin{lemma}[Reference rows with the original pins]
\label{lem:retained-reference-row-pins}
Let \(K\) be the original totally complex number field,
\(\sigma:K\hookrightarrow\mathbb C\) its prescribed embedding,
and \(\epsilon\) the prescribed identity or complex conjugation.
Put \(\sigma_\epsilon=\epsilon\circ\sigma\).
Write \(V_K=K^n\), \(W_K=K^d\) for the original point and
row coordinates.  Let \(E\) be the original finite-dimensional
complex row space, \(j_E\) its \(\sigma\)-semilinear algebraic
embedding, and \(e_{RR}\) the prescribed \(K\)-linear
coordinate isomorphism onto the algebraic row space.  Assume that
the vectors \(j_Ee_{RR}(e_a)\) span \(E\).

In the prescribed complex theta model define the complex linear map
\[
 i_C(z)=\sum_{a=1}^d z_a\,
            \Theta_{\mu_0}(j_R\operatorname{cr}(e_a)).
                                                        \tag{RE1}
\]
Here the composite \(\Theta_{\mu_0}j_R\) is
\(\sigma_\epsilon\)-semilinear, \(\operatorname{cr}\) and
\(e_V\) are \(K\)-linear, and all these maps are the original
ones.  Suppose the
actual coordinate and reference-tuple pins read
\[
\begin{split}
 \iota_{\mu_0}c_F(\sigma y)&=j_Ee_{RR}(P_qy),\\
 e_R(j_Ee_{RR}(w))&=i_C(\sigma_\epsilon w),\\
 \Theta_{\mu_0}(j_Re_Vy)
      &=\Theta_{\mu_0}(j_R\operatorname{cr}(P_qy)),
 \qquad y\in V_K,\quad w\in W_K .
\end{split}                                             \tag{RE2}
\]
The map \(e_R\) is \(\epsilon\)-semilinear.  In the first
line, the original reference row matrix is the identity; thus
this line is the source's \(\mathrm{IotaRd}\) pin after
\(\mathrm{rowMat}_{\mu_0}=I\), not a new choice of coordinates.
Then the reference theta identity has exactly the scalar one:
\[
 \Theta_{\mu_0}(j_Re_Vy)
       =e_R\bigl(\iota_{\mu_0}c_F(\sigma y)\bigr).
                                                        \tag{RE3}
\]
If, in addition, the original good row space is
\(H_\Theta^0=\operatorname{range}i_C\), and its original
lattice norm tie is
\[
 \|i_C(\sigma_\epsilon w)\|^2
       =n_P\|j_Ee_{RR}(w)\|^2\quad(w\in W_K),
       \qquad n_P>0,                                   \tag{RE4}
\]
then \(e_R\) is onto \(H_\Theta^0\) and
\(\|e_RF\|=\sqrt{n_P}\|F\|\) for every \(F\in E\).
The number \(n_P\) is the original row mass.
\end{lemma}

\begin{proof}
Expanding \(w\) in the original coordinate basis gives
\[
 i_C(\sigma_\epsilon w)
       =\Theta_{\mu_0}(j_R\operatorname{cr}(w)).
\]
Substitution into the three equalities (RE2) proves (RE3);
no nonzero comparison scalar is introduced in this calculation.

For the range and norm assertions put
\(g(z)=\sum_a z_a j_Ee_{RR}(e_a)\).
This is onto \(E\), and the second pin in (RE2), first on
the coordinate basis and then by semilinearity, gives
\[
            e_Rg(z)=i_C(\epsilon z).
\]
Coordinatewise \(\epsilon\) is bijective, so the ranges of
\(e_R\) and \(i_C\) coincide.  Since \(K\) is totally
complex, \(\sigma(K)\) contains a nonreal number.  Its
\(\mathbb Q\)-span together with \(1\) is dense in
\(\mathbb C\); hence \(\sigma(K)^d\) is dense in
\(\mathbb C^d\).  Equation (RE4), evaluated at the points
\(z=\sigma w\), and continuity of the two finite-dimensional
norms imply \(\|e_Rg(z)\|^2=n_P\|g(z)\|^2\) for every
\(z\).  Surjectivity of \(g\) gives the assertion.
\end{proof}

\begin{lemma}[The prescribed evaluation fixes the row normalization]
\label{lem:retained-evaluation-normalization}
Let \(\mathcal H_M\) be the original Petersson Hilbert model,
and let \(h:\mathcal H_M\to M\) be its injective realization
in the automorphic multiplicity space, preserving its original
Petersson inner product.  Let \(M^0\subset\operatorname{range}h\)
be the prescribed good multiplicity space.  Suppose the actual
linear row map \(e_E:E\to\mathcal H_M\) satisfies
\[
 \operatorname{range}e_E=h^{-1}(M^0),\qquad
 \|e_EF\|^2=n_P\|F\|^2.                               \tag{RE5}
\]
For the same normalized lowest vector, group identification and
finite placement as in (RP2), assume that the actual evaluation
\(E_t:M^0\to H^0\) is onto the prescribed good automorphic
row space and satisfies
\[
 \|E_tT\|^2=c_{\rm ev}\|T\|_M^2,\qquad c_{\rm ev}>0.
                                                        \tag{RE6}
\]
Then the actual row isometry is
\[
 jF=(c_{\rm ev}n_P)^{-1/2}E_t(h(e_EF)).                 \tag{RE7}
\]
It maps \(E\) onto \(H^0\).  In particular its almost
everywhere formula uses the same \(E_t\) as the finite projector
identity (RP4), with the displayed positive factor.
\end{lemma}

\begin{proof}
The positive norm identity in (RE5) makes \(e_E\) injective.
Together with injectivity of \(h\), its range identity and
\(M^0\subset\operatorname{range}h\), this shows that
\(h e_E:E\to M^0\) is a linear isomorphism.
The norm of its evaluated value is exactly
\[
 \|E_t(h(e_EF))\|^2=c_{\rm ev}n_P\|F\|^2.
\]
Equation (RE6) also makes evaluation injective.  The assumed
onto statement makes the composite onto \(H^0\).  Multiplying
by the displayed positive square root proves (RE7).
If \(E\ne0\), this is the unique positive normalization of
that particular composite which is an isometry.  If \(E=0\),
the formula still defines the same zero isomorphism, without
asserting uniqueness of a scalar from its action on zero vectors.
\end{proof}

Here \(c_{\rm ev}\) is not silently set to one.  In the
source realization it is fixed by the actual Haar comparison
\[
       (\iota^{-1})_*\mu_{A,\mathrm{actual}}
                   =c_{\rm ev}\nu_0 .                 \tag{RE8}
\]
With the original fundamental domains and normalized lowest
vector, quotient integration and the allowed finite translation
give (RE6) with precisely this scalar.  The original
Petersson model, quotient transfer and onto evaluation are
required for this assertion; it is not a consequence of a
finite idempotent alone.

These lemmas refine the application boundary in (RP6)--(RP7).
Once the actual complex-model coordinate, good-range and lattice
norm pins have been supplied, (RP6) may be taken with \(\lambda=1\).
Once the same actual Petersson model, (RE5) and onto evaluation
have been supplied, the row factor is exactly (RE7).
The finite TC/ColIn period dictionary supplies neither (RE4)--(RE6)
for that global model nor the anchored see-saw identity for all
words.  The joined scalar remains
\(a_u=\epsilon(d.a_u)\kappa_u\), and comparison with
\(P_u/\sqrt{c_{0X,u}}\) still retains
\(a_u\sqrt{c_{0X,u}}\).  In particular, these exact row
identities do not establish that its modulus is one, and do not
by themselves transfer an unweighted period trace estimate to
the actual retained theta columns.


\paragraph{Why separate norm identities do not supply this map.}
There is a finite-dimensional obstruction to replacing the common
row identity by the separate norms of the columns. Take two
one-dimensional source spaces and target \(\mathbb C^2\).
Let the two period columns be \(P_1z=ze_1\), \(P_2z=ze_2\),
and let the two candidate theta columns be
\(B_1z=ze_1\), \(B_2z=ze_1\). All four columns are isometries.
Their block-diagonal source Gram operators are both \(I_2\).
No common target
unitary and separate nonzero source scalars can identify the
families, since one pair of image lines is orthogonal and the
other pair coincides. At the same prescribed normalizer \(c=2\),
the target operators are respectively \(I_2\) and
\(\operatorname{diag}(2,0)\); their normalized target excesses
are \(0\) and \(1/2\), although both source deficits vanish.
This example does not refute the actual theta comparison. It
shows why the common row realization is a necessary mathematical
step before transferring the positive-trace estimates.

\begin{lemma}[Restrictions under positive operator majorants]
\label{lem:retained-positive-majorants}
For maps \(T_i:V_i\to W\), put
\[
 S=\bigoplus_iT_i^*T_i,\qquad Q=\sum_iT_iT_i^*,\qquad
 R=\sum_i\dim V_i>0,\qquad D=\dim W>0.
\]
Suppose \(S\ge\alpha I-E_s\) and \(Q\le\beta I+E_t\),
where \(\alpha,\beta\ge0\) and \(E_s,E_t\ge0\).
Restrict by source isometries \(j_i\) and a target projection
\(\Pi\) of rank \(d>0\), with total source rank \(R_0>0\).
For \(c_0>0\), let \(g_0,b_0\) be the prescribed positive
traces of these restricted maps. Then
\[
 \begin{split}
 c_0(g_0+b_0)\le{}&
 \operatorname{tr}(j^*E_sj)+\operatorname{tr}E_t+\beta(D-d)\\
 &+(c_0-\alpha R_0)_++(\beta d-c_0)_+,
 \end{split}
 \tag{RCT4}
\]
where \(j=\bigoplus_i j_i\).
\end{lemma}
\begin{proof}
Write \(Q'=\sum_iT_ij_ij_i^*T_i^*\le Q\). The positive
source loss caused by \(\Pi\) has trace
\(L=\operatorname{tr}((1-\Pi)Q')\), at most
\(\beta(D-d)+\operatorname{tr}((1-\Pi)E_t)\).
The source positive trace is at most
\[
 \operatorname{tr}(j^*E_sj)+L+(c_0-\alpha R_0)_+.
\]
The target positive trace is at most
\(\operatorname{tr}(\Pi E_t\Pi)+(\beta d-c_0)_+\).
These assertions follow from the variational formula for
the trace of a positive part. Adding them proves (RCT4):
the two complementary traces of \(E_t\) sum to its full
trace without requiring any operators to commute.
\end{proof}

Taking \(\alpha=c/R,\ \beta=c/D\) and the original positive
deficit and excess operators gives
\eqref{shared-prescribed-restriction-theta}. In particular
the target excess is charged once under both restrictions.
The formula remains available when the final scale \(c_0\)
is prescribed independently of the retained ranks.

\begin{theorem}[Second moments at the original normalizer]
\label{thm:retained-trace-fixed-moment}
For the actual final maps, let \(S,Q,R,D\) be as above and
\(M=\operatorname{tr}S=\operatorname{tr}Q\). At any
prescribed \(c>0\),
\[
 g_c+b_c=
 \frac{\|S-(c/R)I\|_1+\|Q-(c/D)I\|_1}{2c}.
 \tag{RCT5}
\]
Consequently
\[
 g_c+b_c\le
 \frac{\sqrt{R\operatorname{tr}(S^2)-2cM+c^2}
       +\sqrt{D\operatorname{tr}(Q^2)-2cM+c^2}}{2c}.
 \tag{RCT6}
\]
If known \(\varepsilon,\eta\ge0\) satisfy
\[
 R\|S-(c/R)I\|_{\rm HS}^2\le\varepsilon c^2,\qquad
 D\operatorname{tr}(Q^2)\le(1+\eta)M^2,
\]
then
\[
 g_c+b_c\le
 \frac{\sqrt{\varepsilon}
       +\sqrt{\varepsilon+\eta(1+\sqrt{\varepsilon})^2}}2.
 \tag{RCT7}
\]
These statements include \(M=0\).
\end{theorem}
\begin{proof}
For a Hermitian operator \(A\),
\(\operatorname{tr}A_+=(\|A\|_1+\operatorname{tr}A)/2\).
Apply this first to \((c/R)I-S\), then to \(Q-(c/D)I\).
Their traces \(c-M\) and \(M-c\) cancel, giving (RCT5).
The trace-norm/Hilbert--Schmidt inequality in dimensions
\(R,D\), followed by expansion of the two squares, gives
(RCT6). The two radicands are respectively
\(R\|S-(c/R)I\|_{\rm HS}^2\) and
\(D\|Q-(c/D)I\|_{\rm HS}^2\), so are nonnegative.

Under the last hypotheses, Cauchy--Schwarz gives
\(|M/c-1|\le\sqrt{\varepsilon}\). Therefore
\[
 \frac{D\operatorname{tr}(Q^2)-2cM+c^2}{c^2}
 \le\eta(M/c)^2+(M/c-1)^2
 \le\eta(1+\sqrt{\varepsilon})^2+\varepsilon.
\]
Substitution in (RCT6) proves (RCT7).
\end{proof}

If these formulas give a bound \(E_u\) for the actual
\(g_u+b_u\) at each required parameter, the established
height comparison yields
\[
 P_{\rm norm}(x)\le\kappa(x)^{-1}\operatorname{avg}_u
                  \min\{H_u,E_u(U_u+H_u)\}.
 \tag{RCT8}
\]
Every weight, map, retained subspace and positive normalizer
in this implication is the prescribed one. It neither
constructs the common arithmetic theta/period family nor
provides the numerical second-moment errors for its moving
windows. Those are separate hypotheses, still required
before (RCT8) supplies the final arithmetic budget.
